\section{Background and Preliminaries}
\label{sec:problem}

The input is a finite snapshot of OCDS releases $D$.  A deterministic pipeline, described in
Section \ref{sec:data-method}, turns $D$ into a knowledge graph $G=(V,E)$ and a query-record
projection $U_G$.  Each row $r \in U_G$ is keyed by a contract-node identifier and carries a
controlled set of fields: stored buyer and supplier names, release year, procurement category and
CPV code, selected dates, a monetary value and currency where available, an additivity flag, and
provenance links.  Throughout this paper, \emph{record} means one row of $U_G$.  The projection
keeps two notions of identity apart.  Graph organisation nodes carry canonical identifiers
produced by entity resolution, while query fields such as \texttt{buyer\_name} hold stored strings
taken from the first-party data, so a set of distinct stored names is a set of strings and we read
it that way in every result below.

Let $q$ be a question, $\mathcal{L}$ the closed typed algebra of Section \ref{sec:algebra}, and
$p \in \mathcal{L}$ a program.  Static validation either assigns a root type,
\begin{equation}
  \Gamma \vdash p : \tau,
  \quad
  \tau\in\{\records,\values,\groups,\numbertype,\valuetype,\booltype,\ranking\},
\end{equation}
or returns a structured type error.  Deterministic evaluation then yields an answer envelope
$E_G(p)=(s,\tau,y)$, with execution status $s$ and denotation $y$ when $s=\textsc{ok}$.  A surface
renderer may verbalise $y$; it may not add a factual value that $y$ does not contain.  Evaluation
items divide into three operational statuses.  An item is \emph{answerable} when a program in
$\mathcal{L}$ expresses the computation and its data preconditions hold in the snapshot;
\emph{unsupported} when the question needs an operation outside $\mathcal{L}$, such as a median,
where an explicit abstention is the correct behaviour; and a \emph{no-result} item when a
well-formed scope has an empty denotation, where a faithful empty or zero answer is correct.  The
full runtime adds ambiguity and insufficient evidence as further states, and keeps all of them
distinct, because each calls for a different explanation to the user.

\begin{table}[t]
  \caption{Guarantee boundary.  Only oracle match is an offline answer-correctness measurement;
  none of the first four rows alone establishes real-world truth.}
  \label{tab:guarantee-boundary}
  \centering
  \small
  \begin{tabular}{p{0.18\linewidth}p{0.72\linewidth}}
    \toprule
    Term & Meaning in this paper \\
    \midrule
    Static validity & The object belongs to the closed grammar and all operator input/output
    types compose; depth and node limits are respected. \\
    Executability & The grounded object can be evaluated by the implemented backend without a
    structural or runtime failure. \\
    Runtime verified & The implemented compile, preflight, execution, postflight, sanity, and
    evidence conditions required for release have passed. \\
    Evidence supported & The answer is associated with the record/evidence identifiers emitted
    by execution under the implemented provenance policy. \\
    Oracle match & The typed denotation equals the frozen benchmark oracle under the stated
    exact, set, Boolean, or numeric comparison rule. \\
    Real-world truth & A claim about procurement activity outside the finite snapshot; this is
    not guaranteed by the system or benchmark. \\
    \bottomrule
  \end{tabular}
\end{table}

Five notions of success are easy to conflate, and Table \ref{tab:guarantee-boundary} keeps them
apart so that one inference stays blocked: a program can be type valid and execute cleanly while
expressing the wrong question.  For an answerable item, strict correctness is
\begin{equation}
  \mathbb{1}_{\mathrm{strict}} =
  \mathbb{1}[\Gamma\vdash p:\tau]\;
  \mathbb{1}[E_G(p).s=\textsc{ok}]\;
  \mathbb{1}[\operatorname{match}_{\tau}(E_G(p).y,y^*)],
  \label{eq:strict-correctness}
\end{equation}
where $y^*$ is the frozen oracle and $\operatorname{match}_{\tau}$ is type strict: sets compare as
sets of strings, rankings preserve order and tie semantics, Booleans require Booleans on both
sides, and numbers use the evaluator's documented tolerance (Appendix \ref{app:algebra}).
Unsupported items score only on an explicit refusal of the operation, and no-result items follow
their declared empty-result convention.  Runtime verification is measured separately from
Equation \ref{eq:strict-correctness}.  Where both labels exist for an item, the informative
diagnostic is their $2\times2$ table, which we report only for runs whose trace provenance is
complete (Appendix \ref{app:full-results}).
